\section{Conclusion}
\label{sec:conclusion}

HMRS is a knowledge-document management system whose retrieval path depends on
managed state. It parses heterogeneous sources, assigns content to collections,
derives descriptions and tags, stores documents and structural parts, exposes
catalogue-first access, and records lifecycle state across document files, the
tree registry and per-library metadata. The segment-level gate adds provenance
and an admission signal; it is not a truth oracle or a ranker. This separation
of managed objects, access views and workflow decisions is the system's central
design claim.

The management campaigns show that this loop is operational but not complete.
The October 2026 arXiv campaign recorded 322 parts, with 79 papers retrievable and
21 not at that snapshot; the separate 400-pages-per-language test measured content-based
assignment accuracy of 0.875 for English and 1.0 for Chinese and Japanese. A
33-part storage sample passed its selected checks, while a rename exposed stale
index state and 16 lifecycle scenarios included failures in task-status and
update checks. The system therefore contributes a testable management loop with
explicit failure visibility, rather than an unqualified claim of automatic
consistency.

The access campaigns provide the next layer of evidence. The frozen route-card
audit produced 72/72 full-card versus 37/72 nominally name-only collection
selection, with an unlogged fallback that prevents causal attribution. In the
same-question R3 content-layer follow-up, the tiered chunk-window arm A0h achieved
51/72 (70.8\%), above PD 45/72, A0f 44/72, A0/A1b 39/72 and A2 25/72; its paired
contrast with A0 was 12--0 ($p=0.00049$, exploratory). These values are
 gold-string containment proxies and the arms vary in more than navigation.
The frozen read-all script found all 72 strings but generated no answers, and
R3 read-all was incomplete. Together, the results support the narrower claim that A0h had the highest
observed exact-string containment among completed R3 arms on this corpus and
question set, not semantic answer accuracy, universal retrieval superiority or
a single-factor causal hierarchy claim.

These bounded results support treating the collection as a managed, auditable
object and motivate its use as an inspectable basis for evidence-traced
retrieval. The immediate deployment implication is to expose collection state
and incomplete paths, retain evidence provenance, and separately evaluate
verifier-based admission and ordering before deployment. The global channel can also
silently skip unreadable library indexes while reporting success; until that
condition is surfaced, a successful cross-library response is not evidence of
complete coverage. The study's contribution is a system-level management
contract and a bounded A0h observation, not a claim of universal retrieval
superiority. Broader evaluation needs semantic support review, independently
authored questions, repeated agent trials with run-to-run variance estimates,
larger matched comparisons with multiplicity control, and crash-recovery tests.
